\chapter*{Abstract}
% \phantomsection
% \addcontentsline{toc}{chapter}{Abstract}

Two-dimensional light-element Xenes provide an effective platform for exploring quantum properties generated by reduced dimensionality, strong covalent bonding, and low atomic mass.
Here, the studies on Graphene have established a research paradigm towards atomically thin quantum materials,
however, the zero band gap restricts semiconducting switching performance and band-edge optical response.
This constraint leads us to light-elements allotropes and compound Xenes.
Their properties depend on the bonding network, the number and stacking of layers, and hydrogen termination of the surface.
Changing these features can open a band gap, reshape the optical response, and induce magnetism, superconductivity, or non-trivial band topology.
In this thesis, we apply first-principles calcuation to explore boron- and beryllium-based two-dimensional systems.
We can glimpse the atomic structure, electronic band structure, density of states, dynamical stability, dielectric response, and derived optical properties.
In addition, we also characterizes magnetism, superconductivity,
and non-trivial topology for selected systems.

Our first study examines van der Waals bilayers formed from graphene, borophene, and two-dimensional boron carbides including \ce{BC3} and \ce{B4C3}.
The isolated monolayers demonstrate distinct electronic behaviour.
Graphene is a semimetal and borophene is metallic.
Monolayer \ce{BC3} and \ce{B4C3} are indirect-gap semiconductors,
with their HSE06 band gaps of \qty{1.822}{\electronvolt} and \qty{2.381}{\electronvolt}, respectively.
Through the calculationon of tacking,
we reveal Graphene-Borophene stays metallic,
while Graphene-\ce{BC3} and Graphene-\ce{B4C3} turn semimetallic.
The linear bands from graphene stay near the $\mathrm{K}$ point and open small gaps of \qty{20}{\milli\electronvolt} to \qty{50}{\milli\electronvolt}.
From the dielectric response, we find visible-region absorption and plasmonic excitations in all three heterostructures.
The response differs strongly between in-plane and out-of-plane polarizations.
Graphene-Borophene gives the strongest absorption near \qty{1}{\electronvolt} and shows low-energy plasmonic excitations.
Graphene-\ce{B4C3} shows a similar feature between \qty{5}{\electronvolt} and \qty{6}{\electronvolt}, where its in-plane real dielectric function turns negative.
Below \qty{1}{\electronvolt}, both semimetallic heterostructures also show sharp in-plane energy-loss peaks.

The penta-bipyramid boron study examines boron phases derived from bulk orthorhombic \ce{o-B14}.
Since the boron phases built from \ce{B12} icosahedra, 
\ce{o-B14} consists of edge-sharing \ce{B7} pentagonal bipyramids, and its two-dimensional derivatives retain this building block.
The calculations investigate bulk \ce{o-B14}, monolayer \ce{o-B14}, bilayer \ce{o-B14}, and the hydrogen-terminated bilayer.
The two-dimensional forms display diverse electronic behaviour.
The bilayer is semiconducting with an HSE06 band gap of \qty{0.67}{\electronvolt}.
The monolayer shows magnetism, while the calculations predict phonon-mediated superconductivity in the hydrogen-terminated bilayer.
The optical response is strongly anisotropic.
Several systems absorb in the visible region, and the metallic phases support directional plasmonic modes.
These results demonstrate that the electron-deficient bonding of boron can support semiconductivity, magnetism, superconductivity, and anisotropic optical response in closely related structures.

The beryllene study focuses on pristine beryllene and its hydrogen-functionalized derivatives.
Beryllium represents an extreme light-element limit for elemental two-dimensional materials.
This thesis considers pristine $\alpha$-beryllene, $\beta$-beryllene, trilayer cubic beryllene, which denotes ABA-stacked trilayer beryllene with a square lattice, and hydrogen-functionalized $\alpha$- and $\beta$-beryllene.
The stability analysis confirms the stability of the previously proposed $\alpha$- and $\beta$-beryllene phases and identifies a dynamically stable trilayer cubic beryllene phase.
Hydrogen functionalization strongly modifies the electronic structure.
It drives $\alpha$-beryllene into a wide-gap semiconducting state, whereas the hydrogen-functionalized $\beta$-beryllene structures retain metallic behaviour.
The beryllene optical spectra show strong anisotropy and direction-dependent redistribution of the effective interband and energy-loss response.
We apply parity analysis to the spin-orbit-coupled electronic states and identify conditional non-trivial $\mathbb{Z}_2$ band subspaces in metallic $\alpha$-beryllene and trilayer cubic beryllene.
The smallest sampled direct gaps are of order \qty{1}{\milli\electronvolt} and lie well above the Fermi level.
The selected subspaces of pristine $\beta$-beryllene and the hydrogen-functionalized structures give trivial indices.
Electron-phonon calculations further predict phonon-mediated superconductivity in trilayer cubic beryllene and the hydrogen-functionalized $\beta$-beryllene structures.
The anisotropic Migdal--Eliashberg equations give the highest transition temperature, approximately \qty{31.7}{\kelvin}, for double-sided hydrogen-terminated $\beta$-beryllene relaxed with the D3 dispersion correction.

Our studies demonstrate that the light-element Xenes present complex quantum and optical behaviour via the atomic restructuring in low-dimensional forms.
Heterostructure formation, the bonding network, reduced dimensionality,
and hydrogen functionalization serve as practical routes for tuning the electronic structure, dielectric response, plasmonic modes, superconductivity, magnetism, and band topology.
This thesis provide specific predictions for boron and beryllium allotropes
and outlines a broader strategy for designing ultrathin quantum materials from light elements.
